\documentclass[10pt,a4paper]{amsart}
\usepackage[margin=19mm]{geometry}
\usepackage{amsmath,amssymb,mathtools}
\usepackage[scr=rsfs,cal=euler]{mathalfa}
\usepackage{xfrac}
\usepackage[unicode,hidelinks,bookmarks=false]{hyperref}
\newcommand{\ie}{i.e.\ }
\newcommand{\cf}{cf.\ }
\newcommand{\resp}{respectively\ }
\newcommand{\Subsection}{\subsection}
\providecommand{\nobreakdash}{\nobreak\hbox{-}}
\setlength{\emergencystretch}{2em}
\multlinegap0pt
\usepackage{bm}
\usepackage{bbm}
\usepackage{dsfont}
\usepackage{xspace}
\usepackage{cancel}
\usepackage{mathtools}
\usepackage{amsthm}
\makeatletter
\@ifundefined{theorem}{\newtheorem{theorem}{Theorem}}{}
\makeatother
\newcommand{\TT}{\mathbb{T}}
\title[On Spectra of $\mathbb{T}$-Gain Digraphs]{On Spectra of $\mathbb{T}$-Gain Digraphs}
\author{Shivani Tushar Parab, Tarkeshwar Singh, Mukesh Kumar Nagar}
\date{Source: arXiv:2608.23655v1 (CC BY 4.0)}
\begin{document}
\maketitle

\noindent\emph{Source and licence.} On Spectra of $\mathbb{T}$-Gain Digraphs. Version arXiv:2608.23655v1 (CC BY 4.0), distributed under the Creative Commons Attribution 4.0 International licence, as recorded in the arXiv licence metadata for this exact version and shown on the abstract page https://arxiv.org/abs/2608.23655v1. Full rights notice: licenses/arxiv-2608.23655-CC-BY-4.0.txt.\par\smallskip

\noindent\textit{Editorial scope.} Counted unit: the Theorem whose source label is \texttt{None} in the article (source0.tex lines 589--602, source label \texttt{None}), delivered together with the article's own complete original proof of it (source0.tex lines 604--639, one \texttt{proof} environment of 36 lines). No mathematics is written, repaired or re-derived: statement and proof are the article's own text line for line. Required context retained uncounted: thm:cycle-spectrum (lines 549--551). Every definition, display and label of those lines is retained unchanged. Omitted from this package and not redistributed: 1--548, 552--588, 603--603, 640--920. Nothing inside a retained excerpt is omitted or altered: every retained body is delivered exactly as the article's lines are written, so no omission receipt of any kind accompanies this package. Citations: the retained excerpts carry no citation command at all, so no bibliography entry of the article is transcribed and no reference is redistributed. Reference adaptation: none was needed and none was made; every reference inside the delivered unit resolves inside it, so no unresolved reference or citation is produced. Printed slips kept literal: no printed word, symbol, hypothesis or number of the counted statement or of its proof was altered. Layout and class: the article's own class is \texttt{article} with packages including amsmath, amsthm, amssymb, geometry, hyperref, microtype, tikz, enumitem, times; the wrapper is the lane's pinned \texttt{amsart} editorial wrapper at 10pt on A4, which restates the theorem-like environments the retained text needs and adds the article's own macro definitions that the retained text needs (\texttt{\textbackslash TT}), with extra wrapper packages (bm, bbm, dsfont, xspace, cancel, mathtools). The delivered numbering is the unit's own. Assets: no retained span contains a figure environment, a graphics-inclusion command, a PDF-page inclusion, a table or a TikZ picture; no image file is copied into this package, the package declares no asset dependency and no image file is redistributed with it. Licence: version arXiv:2608.23655v1 of this article is distributed under the Creative Commons Attribution 4.0 International licence, as recorded in the arXiv licence metadata for this exact version and shown on the abstract page https://arxiv.org/abs/2608.23655v1. A full rights notice is delivered at licenses/arxiv-2608.23655-CC-BY-4.0.txt. Characters: every retained excerpt is pure ASCII, so the delivered engine, XeLaTeX with its default Latin Modern text font, reports no missing character.
\begin{theorem}\label{thm:cycle-spectrum}
Let $\Phi=(C_n,\varphi)$ be a $\mathbb{T}$-gain digraph on the directed cycle $v_1\to v_2\to\cdots\to v_n\to v_1$. If $g=\varphi(C_n)=\prod_{i=1}^{n}\varphi(v_i,v_{i+1})$, then $\operatorname{spec}(A(\Phi)) = \{\lambda\in\mathbb{C}:\lambda^n=g\}$ and $\rho(\Phi)=1$.
\end{theorem}

\begin{theorem}
Let $\Phi=(D,\varphi)$ be a $\TT$-gain unicyclic digraph on $n$ vertices whose unique directed cycle is $C_\ell$, and let
$g=\varphi(C_\ell)\in\TT$.
Then
\[
\operatorname{spec}(A(\Phi))
=
\{0^{\,n-\ell}\}
\cup
\{\lambda\in\mathbb C:\lambda^\ell=g\},
\]
where $0^{\,n-\ell}$ denotes the eigenvalue $0$ with multiplicity
$n-\ell$. Consequently, $\rho(\Phi)=1.$
\end{theorem}

\begin{proof}
Order the vertices which are not lying on the directed cycle first (from the tips toward the cycle), followed by the vertices of the directed cycle.
With this ordering,
\[
A(\Phi)=
\begin{pmatrix}
N&R\\
0&C(\Phi)
\end{pmatrix},
\]
where $N$ is a strictly upper triangular matrix corresponding to the arcs outside the directed circle, $R$ contains the arcs joining the tails to the cycle, and $C(\Phi)$ is the gain adjacency matrix of the directed cycle
$C_\ell$.

Since $N$ is strictly upper triangular, it is nilpotent. Hence all
eigenvalues of $N$ are equal to $0$, with algebraic multiplicity
$n-\ell$. Since $A(\Phi)$ is block upper triangular, its spectrum is the union of the spectra of its diagonal blocks. Therefore,
$\operatorname{spec}(A(\Phi))
=
\operatorname{spec}(N)
\cup
\operatorname{spec}(C(\Phi))$.
By Theorem~\ref{thm:cycle-spectrum},
$\operatorname{spec}(C(\Phi))
=
\{\lambda\in\mathbb C:\lambda^\ell=g\}.$
Thus
\[
\operatorname{spec}(A(\Phi))
=
\{0^{\,n-\ell}\}
\cup
\{\lambda\in\mathbb C:\lambda^\ell=g\}.
\]
Finally, since $g\in\TT$, every solution of $\lambda^\ell=g$ has modulus one. Hence every nonzero eigenvalue of $A(\Phi)$ has modulus one, while the remaining eigenvalues are zero. Therefore,
$\rho(\Phi)=1.$
\end{proof}

\end{document}
